\par\addvspace{12pt}\Needspace{9\baselineskip}

\section{综合习题}

\label{sol:sec:analysis:9}

% ITEM: analysis-difference-to-sum
\solutionitem{analysis-difference-to-sum}{习题1.4（1）}
设\(a_1=1,a_{n+1}-a_n=2n\)。

\quad \exercisesubitem{1}用累加法求通项，并用数学归纳法核对。

\quad \exercisesubitem{2}求前\(n\)项和\(S_n\)，再由\(S_n-S_{n-1}\)恢复\(a_n\)。
\solutionanswer
（1）将相邻差累加，\(a_n=a_1+\sum_{k=1}^{n-1}2k=1+n(n-1)=n^2-n+1\)。核对：\(n=1\)成立；若\(a_n=n^2-n+1\)，则\(a_{n+1}=a_n+2n=n^2+n+1=(n+1)^2-(n+1)+1\)，归纳得证。（2）\(S_n=\sum_{k=1}^n(k^2-k+1)=n(n+1)(2n+1)/6-n(n+1)/2+n=n(n^2+2)/3\)。因此\(S_n-S_{n-1}=n^2-n+1=a_n\)，亦恢复原通项。


% ITEM: analysis-arithmetic-telescoping
\solutionitem{analysis-arithmetic-telescoping}{习题1.4（2）}
等差数列的前\(n\)项和为\(S_n\)，且\(S_3=12,S_6=42\)。求通项，并求
\begin{equation*}
 T_n=\sum_{k=1}^n\frac1{a_ka_{k+1}}
\end{equation*}
及\(\lim_{n\to\infty}T_n\)。
\solutionanswer
设首项\(a\)、公差\(d\)。由\(3a+3d=12\)及\(6a+15d=42\)，得\(a=2,d=2\)，故\(a_k=2k\)。于是
\[T_n=\sum_{k=1}^n\frac1{(2k)(2k+2)}=\frac14\sum_{k=1}^n\left(\frac1k-\frac1{k+1}\right)=\frac14\left(1-\frac1{n+1}\right),\quad \lim T_n=\frac14.\]


% ITEM: analysis-negative-geometric
\solutionitem{analysis-negative-geometric}{习题1.4（3）}
等比数列的首项为\(3\)，公比为\(-1/2\)。求通项、前\(n\)项和及二者的极限。原数列与部分和数列是否单调？它们的奇、偶子列分别怎样变化？
\solutionanswer
\(a_n=3(-1/2)^{n-1}\)，\(S_n=2(1-(-1/2)^n)\)。故\(a_n\to0,S_n\to2\)。原数列符号交替，不单调；奇数子列\(a_{2k-1}=3/4^{k-1}\)递减至0，偶数子列\(a_{2k}=-3/2^{2k-1}\)递增至0。部分和\(S_n\)奇数项大于2且递减至2，偶数项小于2且递增至2，整体不单调。


% ITEM: analysis-odd-weighted-sum
\solutionitem{analysis-odd-weighted-sum}{习题1.4（4）}
设\(S_n=\sum_{k=1}^n(2k-1)3^{k-1}\)。

\quad \exercisesubitem{1}用错位相减求\(S_n\)。

\quad \exercisesubitem{2}分别求\(S_n/3^n\)、\(S_n/(n3^n)\)的极限，比较两种尺度下的结果。
\solutionanswer
令\(S_n=\sum_{k=1}^n(2k-1)3^{k-1}\)。错位相减得
\[3S_n-S_n=\sum_{k=1}^n(2k-1)3^k-\sum_{k=1}^n(2k-1)3^{k-1}=(2n-1)3^n-1-2\sum_{k=1}^{n-1}3^k,\]
从而\(S_n=((2n-1)3^n-3^n+2)/2=((2n-2)3^n+2)/2=(n-1)3^n+1\)。因此\(S_n/3^n=n-1+3^{-n}\to+\infty\)，而\(S_n/(n3^n)=1-1/n+1/(n3^n)\to1\)。前一尺度未消去随\(n\)增长的权重，故发散；除以额外的\(n\)后得到有限极限。


% ITEM: analysis-radical-sum
\solutionitem{analysis-radical-sum}{习题1.4（5）}
设
\begin{equation*}
 S_n=\sum_{k=1}^n\frac1{\sqrt{k+1}+\sqrt k}.
\end{equation*}
求\(S_n\)，并分别求\(\lim_{n\to\infty}S_n/n\)、\(\lim_{n\to\infty}S_n/\sqrt n\)。说明同一个部分和数列在不同尺度下的极限有何不同。
\solutionanswer
有\(1/(\sqrt{k+1}+\sqrt k)=\sqrt{k+1}-\sqrt k\)，故\(S_n=\sqrt{n+1}-1\)。于是\(S_n/n\to0\)，而\(S_n/\sqrt n=\sqrt{1+1/n}-1/\sqrt n\to1\)。同一数列相对线性尺度趋零，相对平方根尺度趋于非零有限值1。


% ITEM: analysis-epsilon-order
\solutionitem{analysis-epsilon-order}{习题1.4（6）}
\exercisesubitem{1}只用定义证明\((3n-1)/(n+2)\to3\)，并对给定的\(\varepsilon>0\)写出一个可用的\(N\)。

\quad \exercisesubitem{2}设\(a_n\to a,b_n\to b\)，且从某项起\(a_n\le b_n+1/n\)，证明\(a\le b\)。
\solutionanswer
（1）\[\left|\frac{3n-1}{n+2}-3\right|=\frac7{n+2}.\]给定\(\varepsilon>0\)，取正整数\(N>7/\varepsilon\)，则\(n\ge N\)时误差\(7/(n+2)<\varepsilon\)，按定义极限为3。（2）若\(a_n\to a,b_n\to b\)，且尾部\(a_n\le b_n+1/n\)，取极限保序（或由\(a_n-b_n\le1/n\)及\(a_n-b_n\to a-b\)），得\(a-b\le0\)，即\(a\le b\)。


% ITEM: analysis-cauchy-error
\solutionitem{analysis-cauchy-error}{习题1.4（7）}
设实数列满足\(|a_{n+1}-a_n|\le2^{-n}\)。用Cauchy准则证明它收敛，并给出极限\(a\)与\(a_n\)之间的误差上界。若条件换为\(|a_{n+1}-a_n|\le1/n\)，是否仍必收敛？说明理由。
\solutionanswer
当\(m>n\)，
\[|a_m-a_n|\le\sum_{k=n}^{m-1}|a_{k+1}-a_k|\le\sum_{k=n}^{m-1}2^{-k}<2^{1-n}.\]
右端随\(n\to\infty\)趋零，故列满足Cauchy准则而收敛。令\(m\to\infty\)，得\(|a-a_n|\le2^{1-n}\)。换成\(1/n\)时不能保证：取\(a_n=\sum_{k=1}^{n-1}1/k\)，相邻差恰为\(1/n\)，但调和部分和发散。


% ITEM: analysis-parity-limit
\solutionitem{analysis-parity-limit}{习题1.4（8）}
若数列\(\{a_n\}\)的奇子列与偶子列都收敛于\(a\)，证明原数列也收敛于\(a\)。
\solutionanswer
给定\(\varepsilon>0\)，奇子列、偶子列均趋于\(a\)，故存在指标阈值，使所有足够后的奇数项及偶数项分别与\(a\)相差小于\(\varepsilon\)。取两阈值较大者，则每个足够大的正整数不是奇即偶，故\(|a_n-a|<\varepsilon\)，原数列趋于\(a\)。


% ITEM: analysis-bisection
\solutionitem{analysis-bisection}{习题1.4（9）}
从\([a_0,b_0]=[1,2]\)开始，每次取\(m_n=(a_n+b_n)/2\)。若\(m_n^3\le2\)，保留\([m_n,b_n]\)；否则保留\([a_n,m_n]\)，记为\([a_{n+1},b_{n+1}]\)。

\quad \exercisesubitem{1}证明这些区间构成长度趋于零的区间套，公共点\(\xi\)满足\(\xi^3=2\)。

\quad \exercisesubitem{2}用进行十次二分后的区间中点近似\(\xi\)，给出误差上界。这里的收敛性为何不需要先求出\(\xi\)？
\solutionanswer
每次二分都保持\(a_n^3\le2\le b_n^3\)，区间长度又满足\(b_n-a_n=2^{-n}\to0\)。由闭区间套定理，两端点都趋于唯一公共点\(\xi\)。利用立方的极限法则，由\(a_n^3\le2\le b_n^3\)得到\(\xi^3\le2\le\xi^3\)，故\(\xi^3=2\)。十次后按规则逐次二分得
\[[a_{10},b_{10}]=[1.259765625,1.2607421875],\quad m_{10}=1.26025390625.\]
误差不超过半区间长\((b_{10}-a_{10})/2=1/2048\approx0.00048828125\)。此构造的区间端点及长度已能证明中点逼近其公共点；无需预先写出\(\xi\)的显式值。


% ITEM: analysis-stolz-weighted
\solutionitem{analysis-stolz-weighted}{习题1.4（10）}
利用适当的极限工具，计算或证明：
\begin{equation*}
 \lim_{n\to\infty}\frac{\sqrt1+\sqrt2+\cdots+\sqrt n}{n^{3/2}}=\frac23;
 \qquad
 x_n\to l\ \Longrightarrow\
 \frac{\sum_{k=1}^n kx_k}{n^2}\to\frac l2.
\end{equation*}
说明相邻增量的比较怎样确定上述极限。
\solutionanswer
设\(A_n=\sum_{k=1}^n\sqrt k\)。分母\(n^{3/2}\)严格递增且趋于无穷，适合使用Stolz定理。由立方差和平方差恒等式，
\[
(n+1)^{3/2}-n^{3/2}
=\frac{2n+1+\sqrt{n(n+1)}}{\sqrt{n+1}+\sqrt n}.
\]
因此
\[
\frac{A_{n+1}-A_n}{(n+1)^{3/2}-n^{3/2}}
=\frac{n+1+\sqrt{n(n+1)}}{2n+1+\sqrt{n(n+1)}}\to\frac23.
\]
由Stolz定理，第一极限为\(2/3\)。
第二式令\(A_n=\sum_{k=1}^n kx_k,B_n=n^2\)。相邻增量比为
\[\frac{A_{n+1}-A_n}{B_{n+1}-B_n}=\frac{(n+1)x_{n+1}}{2n+1}\to\frac l2.\]
由Stolz定理得所求极限\(l/2\)。两处均以相邻增量之比的极限决定累积和与分母的比值。


% ITEM: analysis-e-error
\solutionitem{analysis-e-error}{习题1.4（11）}
记\(E_n=\sum_{k=0}^n1/k!\)。

\quad \exercisesubitem{1}由\(0<e-E_n<1/(n\,n!)\)，求使这个误差上界小于\(10^{-6}\)的最小正整数\(n\)。

\quad \exercisesubitem{2}利用本章关于余项的结论，证明\(n\,n!(e-E_n)\to1\)。误差上界与实际误差的比值趋于多少？
\solutionanswer
（1）误差上界要求\(1/(n n!)<10^{-6}\)，即\(n n!>10^6\)。\(9\cdot9!=3,265,920>10^6\)，而\(8\cdot8!=322,560<10^6\)，故最小正整数\(n=9\)。（2）余项
\[e-E_n=\sum_{j=1}^{\infty}\frac1{(n+j)!}=\frac1{n!}\sum_{j=1}^{\infty}\frac1{(n+1)(n+2)\cdots(n+j)}.\]
乘以\(n n!\)，首项为\(n/(n+1)\to1\)，其余项和趋于0：对\(j\ge2\)，乘积至少为\((n+1)^j\)，故余项总和至多\(n\sum_{j=2}^\infty(n+1)^{-j}=1/n\)。所以\(n n!(e-E_n)\to1\)。给定上界为\(1/(n n!)\)，其与实际误差之比为\([1/(n n!)]/(e-E_n)\to1\)。


% ITEM: analysis-newton-square-root
\solutionitem{analysis-newton-square-root}{习题1.4（12）}
设\(A>0,x_1>0\)，并令
\begin{equation*}
 x_{n+1}=\frac12\left(x_n+\frac A{x_n}\right).
\end{equation*}
证明从第二项起\(x_n\ge\sqrt A\)，且数列单调不增，进而证明\(x_n\to\sqrt A\)。再证明误差关系
\begin{equation*}
 x_{n+1}-\sqrt A=\frac{(x_n-\sqrt A)^2}{2x_n},
\end{equation*}
说明它为何适合用于平方根的近似计算。
\solutionanswer
记\(r=\sqrt A\)。递推保持正性，由均值不等式，\(x_{n+1}=(x_n+A/x_n)/2\ge r\)。因此从第二项起，
\[
x_{n+1}-x_n=\frac{A-x_n^2}{2x_n}\le0.
\]
尾列单调不增且以下界\(r\)控制，设极限为\(L\ge r\)。取极限得\(L=(L+A/L)/2\)，故\(L=r\)。再通分得到
\[
x_{n+1}-r=\frac{x_n^2-2rx_n+r^2}{2x_n}
=\frac{(x_n-r)^2}{2x_n}.
\]
在\(x_n\ge r\)时，新误差不超过旧误差平方的\(1/(2r)\)倍；靠近极限后，平方关系使误差迅速减小。


% ITEM: analysis-asymptotic-radicals
\solutionitem{analysis-asymptotic-radicals}{习题1.4（13）}
计算下列极限，并说明怎样消去式中的主要部分：
\begin{equation*}
 \exercisesubitem{1}\ \lim_{n\to\infty}\left(\sqrt{n^2+n}-n\right);
 \qquad
 \exercisesubitem{2}\ \lim_{n\to\infty}n\left(\sqrt{n^2+n}-n-\frac12\right).
\end{equation*}
\solutionanswer
（1）有理化消去两个同阶主项：
\[
\sqrt{n^2+n}-n=\frac n{\sqrt{n^2+n}+n}\to\frac12.
\]
（2）再减去常数主项，利用\((n+1/2)^2=n^2+n+1/4\)：
\[
n\left(\sqrt{n^2+n}-n-\frac12\right)
=\frac{-n/4}{\sqrt{n^2+n}+n+1/2}\to-\frac18.
\]
两步都由平方差恒等式控制消去后的剩余量，无须使用展开公式。


% ITEM: analysis-geometric-averages
\solutionitem{analysis-geometric-averages}{习题1.4（14）}
设\(a_n\to l\)。证明
\begin{equation*}
 \lim_{n\to\infty}\frac{a_1+2a_2+\cdots+2^{n-1}a_n}{2^n}=l.
\end{equation*}
若将条件改为左侧数列收敛于\(l\)，是否能反推出\(a_n\to l\)？与Cauchy命题的逆命题比较。
\solutionanswer
记\(A_n=2^{-n}\sum_{k=1}^n2^{k-1}a_k\)。权重之和为\(1-2^{-n}\)，故
\[
A_n-l=2^{-n}\sum_{k=1}^n2^{k-1}(a_k-l)-2^{-n}l.
\]
给定\(\varepsilon>0\)，取\(K\)使\(k>K\)时\(|a_k-l|<\varepsilon\)。于是
\[
|A_n-l|\le2^{-n}\sum_{k=1}^K2^{k-1}|a_k-l|+\varepsilon+2^{-n}|l|.
\]
固定前段与最后一项都趋于零，再由\(\varepsilon\)任意即得\(A_n\to l\)。

逆命题也成立。相邻两次加权平均满足\(2A_n-A_{n-1}=a_n\)，故若\(A_n\to l\)，则\(a_n\to2l-l=l\)。Cauchy命题的逆命题一般不成立，例如\((-1)^n\)的算术平均趋于零而原列不收敛；这里最新一项的权重恒为\(1/2\)，能从相邻平均中恢复出来。


% ITEM: analysis-shifted-telescoping
\solutionitem{analysis-shifted-telescoping}{习题1.4（15）}
设\(\alpha>-1\)。求\(\sum_{n=1}^{\infty}1/[(n+\alpha)(n+\alpha+1)]\)的和，并求前\(N\)项和的误差。对给定的\(\varepsilon>0\)，写出保证误差小于\(\varepsilon\)的\(N\)应满足的条件。
\solutionanswer
裂项后有
\[
S_N=\sum_{n=1}^N\left(\frac1{n+\alpha}-\frac1{n+\alpha+1}\right)
=\frac1{1+\alpha}-\frac1{N+1+\alpha}.
\]
级数的和为\(1/(1+\alpha)\)，余项恰为\(1/(N+1+\alpha)\)。保证误差小于\(\varepsilon\)的条件为\(N+1+\alpha>1/\varepsilon\)，据此取正整数\(N\)即可。


% ITEM: analysis-monotone-summable
\solutionitem{analysis-monotone-summable}{习题1.4（16）}
设\(a_n>0\)单调不增，且\(\sum_{n=1}^{\infty}a_n\)收敛。

\quad \exercisesubitem{1}通过比较\(na_{2n}\)与\(\sum_{k=n+1}^{2n}a_k\)，证明\(na_n\to0\)。

\quad \exercisesubitem{2}若不要求单调，结论是否仍成立？试用只在少数下标处较大的数列说明。
\solutionanswer
（1）单调性使一个项能由相邻一整段的和控制：
\[
0<na_{2n}\le\sum_{k=n+1}^{2n}a_k\to0.
\]
这里右边是收敛级数的尾部区段。若下标为\(2n\)，则\(2na_{2n}\to0\)；若下标为\(2n+1\)，由单调性又有
\[
0<(2n+1)a_{2n+1}\le(2n+1)a_{2n}
=\left(2+\frac1n\right)na_{2n}\to0.
\]
因此全列\(na_n\to0\)。

（2）去掉单调性不成立。取
\[
a_n=\frac1{n^2}+\begin{cases}1/j^2,&n=2^j,\ j\ge1,\\0,&\text{其余下标}.\end{cases}
\]
各项严格为正，且总和不超过两个收敛的平方倒数级数之和；但\(2^ja_{2^j}\ge2^j/j^2\to+\infty\)。集中在少数下标上的较大项破坏了结论。


% ITEM: analysis-reciprocal-iteration
\solutionitem{analysis-reciprocal-iteration}{习题1.4（17）}
设\(x_1>0\)，且\(x_{n+1}=x_n/(1+x_n)\)。

\quad \exercisesubitem{1}考察\(1/x_n\)，求\(x_n\)的通项与\(\lim_{n\to\infty}nx_n\)。

\quad \exercisesubitem{2}证明\(x_n-x_{n+1}=x_nx_{n+1}\)，并求级数\(\sum_{n=1}^{\infty}x_nx_{n+1}\)的和与前\(N\)项和的误差。
\solutionanswer
（1）取倒数得\(1/x_{n+1}=1/x_n+1\)，因此
\[
x_n=\frac1{n-1+1/x_1},\qquad nx_n\to1.
\]
（2）由递推式，\(x_n-x_{n+1}=x_n^2/(1+x_n)=x_nx_{n+1}\)。所以前\(N\)项和为\(x_1-x_{N+1}\)，级数的和为\(x_1\)，余项恰为\(x_{N+1}=1/(N+1/x_1)\)。


% ITEM: analysis-unique-cluster
\solutionitem{analysis-unique-cluster}{习题1.4（18）}
设\(\{a_n\}\)有界，且它的每个收敛子列都以\(a\)为极限。证明\(a_n\to a\)。若原数列不收敛于\(a\)，试先找出始终与\(a\)相隔一定距离的子列，再使用有界数列的收敛子列定理。
\solutionanswer
若\(a_n\)不趋于\(a\)，则存在\(\varepsilon_0>0\)及子列\(a_{n_k}\)，使每一项都满足\(|a_{n_k}-a|\ge\varepsilon_0\)。这个子列仍有界，由有界数列的收敛子列定理，可再取一个收敛子列，设其极限为\(b\)。取极限得\(|b-a|\ge\varepsilon_0\)，与题设所有收敛子列都以\(a\)为极限矛盾。故\(a_n\to a\)。


% ITEM: analysis-bounded-sums
\solutionitem{analysis-bounded-sums}{习题1.4（19）}
设\(S_n=\sum_{k=1}^n a_k\)有界。

\quad \exercisesubitem{1}证明\((a_1+\cdots+a_n)/n\to0\)。这是否说明\(a_n\to0\)，或说明\(\sum a_n\)收敛？分别说明理由。

\quad \exercisesubitem{2}若另有\(a_n\ge0\)，结论有什么变化？比较有界性与单调性在这里的作用。
\solutionanswer
（1）由\(|S_n|\le M\)，得\(|S_n/n|\le M/n\to0\)。这不保证\(a_n\to0\)，也不保证\(\sum a_n\)收敛：取\(a_n=(-1)^n\)，部分和有界，但通项不趋于零，部分和也不收敛。

（2）若另有\(a_n\ge0\)，则\(S_n\)单调不减且有界，因而收敛；于是级数收敛，通项\(a_n\to0\)。有界性提供控制，非负性提供单调性，两者共同保证了部分和收敛。


% ITEM: analysis-nested-intervals
\solutionitem{analysis-nested-intervals}{习题1.4（20）}
设闭区间\([u_{n+1},v_{n+1}]\subseteq[u_n,v_n]\)，且\(v_n-u_n\to d\ge0\)。

\quad \exercisesubitem{1}证明两端点数列分别收敛，记极限为\(u,v\)，并证明全部区间的公共部分恰为\([u,v]\)，其长度为\(d\)。

\quad \exercisesubitem{2}若改为开区间，即使长度趋于零，是否仍一定有公共点？举例说明。
\solutionanswer
（1）包含关系给出\(u_n\)单调不减、\(v_n\)单调不增，且\(u_1\le u_n\le v_n\le v_1\)，因此两列分别收敛于\(u,v\)。由长度的极限知\(v-u=d\)。对任意固定的\(n\)，让\(m\ge n\)趋于无穷，得到\(u_n\le u\le v\le v_n\)，所以\([u,v]\)包含于每个区间。反过来，若\(x\)属于所有区间，则由\(u_n\le x\le v_n\)取极限知\(u\le x\le v\)。故交集恰为\([u,v]\)，长度为\(d\)。

（2）开区间不保证有公共点。例如\((0,1/n)\)逐层包含且长度趋于零，交集却为空。


% ITEM: analysis-euler-block
\solutionitem{analysis-euler-block}{习题1.4（21）}
记\(H_n=\sum_{k=1}^n1/k\)，已知\(\gamma_n=H_n-\ln n\)收敛。对固定整数\(m\ge2\)，求\(\lim_{n\to\infty}(H_{mn}-H_n)\)。为什么不能把含有\(n\)项的这个和，直接拆成逐项极限之和？
\solutionanswer
直接利用已经收敛的差额，得到
\[
H_{mn}-H_n=\ln(mn)-\ln n+\gamma_{mn}-\gamma_n\to\ln m.
\]
这里求和项数随\(n\)改变，不能使用固定有限和的逐项极限法则。原和中各项虽越来越小，总和却趋于\(\ln m\)，这正是两种取极限方式不能混同的例子。


% ITEM: analysis-alternating-radicals
\solutionitem{analysis-alternating-radicals}{习题1.5（1）〔选做〕}
判断\(\sum_{n=1}^{\infty}(-1)^{n-1}(\sqrt{n+1}-\sqrt n)\)是否收敛、是否绝对收敛，并给出前\(N\)项和近似级数和的误差上界。若要求误差不超过\(10^{-2}\)，应怎样选择\(N\)？
\solutionanswer
令\(c_n=\sqrt{n+1}-\sqrt n=1/(\sqrt{n+1}+\sqrt n)\)。则\(c_n\)单调递减且趋于零，交错级数收敛；但\(\sum_{n=1}^Nc_n=\sqrt{N+1}-1\to+\infty\)，故不绝对收敛。

由交错余项估计，
\[
|S-S_N|\le c_{N+1}=\frac1{\sqrt{N+2}+\sqrt{N+1}}
<\frac1{2\sqrt{N+1}}.
\]
取\(N\ge2499\)，便能保证误差不超过\(10^{-2}\)。
\par
